\chapter{Trabalhos Relacionados}
\label{cap:trabalhos-relacionados}

Esta seção discute estudos relacionados ao uso de modelos de linguagem de grande porte em comunicação em saúde, letramento em saúde, simplificação textual e mitigação de riscos. A análise busca posicionar este trabalho em relação à literatura existente, destacando aproximações, diferenças e lacunas. Em especial, são considerados estudos que investigam LLMs na educação de pacientes, aplicações mHealth, simplificação de materiais de saúde e uso de guardrails para reduzir riscos associados à geração automática de texto.

\section{LLMs em educação de pacientes e comunicação em saúde}
\label{sec:llms-educacao-pacientes-comunicacao-saude}

O uso de modelos de linguagem de grande porte em educação de pacientes tem sido investigado em diferentes contextos da medicina, especialmente pela capacidade desses modelos de gerar explicações em linguagem natural, responder dúvidas e reformular conteúdos técnicos para públicos não especializados. Aydin et al. ([VERIFICAR\_CITACAO: AYDIN et al., 2024]) apresentam uma revisão de escopo sobre aplicações de LLMs na educação e no engajamento de pacientes. A revisão identificou 201 estudos e organizou os usos encontrados em seis temas: geração de materiais educativos, interpretação de informações médicas, recomendações de estilo de vida, apoio ao uso personalizado de medicamentos, instruções de cuidado perioperatório e melhoria da interação médico-paciente ([VERIFICAR\_CITACAO: AYDIN et al., 2024]).

Esse trabalho é relevante para este TCC porque mostra que LLMs já vêm sendo estudados como ferramentas de apoio à comunicação entre informações médicas e pacientes. Aydin et al. ([VERIFICAR\_CITACAO: AYDIN et al., 2024]) indicam que LLMs podem auxiliar na criação de materiais mais acessíveis e na tradução de informações médicas para linguagem mais amigável ao paciente. Ao mesmo tempo, os autores destacam limitações relacionadas à legibilidade, à acurácia, a vieses e à necessidade de assegurar confiabilidade do conteúdo gerado. Essas limitações são particularmente importantes para este TCC, pois a simplificação textual em saúde não pode ser avaliada apenas pela fluidez e naturalidade da linguagem, mas também pela preservação das informações médicas essenciais.

A principal diferença entre a revisão de Aydin et al. ([VERIFICAR\_CITACAO: AYDIN et al., 2024]) e este trabalho está no nível de especificidade. A revisão discute aplicações amplas de LLMs em educação e engajamento de pacientes, abrangendo desde respostas a perguntas até materiais educacionais e interação médico-paciente. Este TCC, por sua vez, concentra-se em uma tarefa mais delimitada: simplificar e adaptar linguisticamente textos técnicos de planos de cuidado para perfis de pacientes em um ambiente mHealth. Assim, enquanto Aydin et al. ([VERIFICAR\_CITACAO: AYDIN et al., 2024]) ajudam a justificar a relevância geral dos LLMs na comunicação em saúde, este trabalho avança em uma aplicação específica, com foco em textos já existentes, guardrails para delimitação da geração e redução de riscos e avaliação por indicadores de simplificação linguística.

Outro ponto importante é que parte da literatura sobre LLMs em educação de pacientes avalia respostas geradas diretamente pelos modelos a perguntas clínicas ou dúvidas frequentes. Neste TCC, a proposta é mais restrita: o modelo não é usado para produzir aconselhamento clínico novo, mas para reescrever conteúdos previamente definidos ou validados por profissionais de saúde. Essa delimitação reduz o escopo clínico da geração e aproxima a solução de uma ferramenta de apoio comunicacional, não de decisão médica.

\section{LLMs, letramento em saúde e mHealth}
\label{sec:llms-letramento-saude-mhealth}

Entre os trabalhos analisados, Loughran et al. ([VERIFICAR\_CITACAO: LOUGHRAN et al., 2024]) é um dos mais próximos do tema deste TCC por conectar explicitamente LLMs, letramento em saúde e mHealth. O estudo discute o uso de modelos de linguagem como apoio à adaptação de materiais educacionais em saúde para aplicações mHealth, com o objetivo de reduzir barreiras de compreensão associadas à complexidade textual. Os autores avaliam o uso de ChatGPT para reduzir a complexidade linguística de conteúdos já produzidos, tratando o modelo como um recurso de apoio à reescrita e à edição textual, e não apenas como uma ferramenta de geração aberta de respostas.

A proximidade com este TCC está na preocupação em tornar conteúdos de saúde mais compreensíveis no contexto de aplicações digitais. Ambos os trabalhos partem da ideia de que a efetividade de uma aplicação mHealth depende não apenas da disponibilização da informação, mas da capacidade do paciente de compreendê-la adequadamente. A diferença principal está no escopo da implementação e nos mecanismos de controle considerados. Enquanto Loughran et al. ([VERIFICAR\_CITACAO: LOUGHRAN et al., 2024]) investigam o uso de LLMs como apoio à redução de complexidade textual em materiais educacionais, este TCC propõe uma implementação funcional voltada à simplificação e adaptação linguística de textos técnicos de saúde. Além disso, este trabalho enfatiza a adaptação ao perfil do paciente, considerando características como idade, escolaridade e nível de letramento em saúde, e incorpora guardrails para delimitar o uso do modelo, restringir a geração ao conteúdo fornecido e reduzir riscos como criação de informações novas, alterações de sentido ou respostas fora do escopo da simplificação textual.

\section{Simplificação de textos médicos com LLMs}
\label{sec:simplificacao-textos-medicos-llms}

A simplificação de textos médicos com apoio de LLMs tem sido investigada como uma forma de tornar informações técnicas mais compreensíveis para pacientes e usuários sem formação na área da saúde. Um trabalho relevante nessa direção é o de Jeblick et al. ([VERIFICAR\_CITACAO: JEBLICK et al., 2024]), que avalia o uso do ChatGPT para simplificar laudos radiológicos. No estudo, relatórios de radiologia foram simplificados com o modelo e posteriormente avaliados por radiologistas quanto à correção factual, completude e potencial de dano ao paciente. Os resultados indicaram potencial no uso de LLMs para tornar textos médicos mais acessíveis, mas também apontaram a ocorrência de informações incorretas, omissões e trechos que poderiam levar a interpretações inadequadas por parte dos pacientes ([VERIFICAR\_CITACAO: JEBLICK et al., 2024]).

Esse trabalho se aproxima deste TCC por tratar diretamente da simplificação de linguagem médica com LLMs e por reconhecer que a qualidade da simplificação não depende apenas da redução da complexidade textual, mas também da preservação do conteúdo original. No entanto, o estudo de Jeblick et al. ([VERIFICAR\_CITACAO: JEBLICK et al., 2024]) está concentrado em laudos radiológicos, enquanto este TCC trabalha com textos técnicos de saúde voltados à orientação do paciente. Além disso, este trabalho considera a adaptação linguística ao perfil do paciente e incorpora mecanismos de controle para restringir a geração ao conteúdo fornecido e reduzir riscos de alteração de sentido ou criação de informações novas.

Em síntese, o estudo reforça que LLMs podem apoiar a simplificação de textos médicos, mas que essa tarefa exige cuidado para evitar omissões, distorções e interpretações inadequadas. Essa constatação se relaciona diretamente com a proposta deste TCC, que busca combinar simplificação linguística, adaptação ao perfil do paciente, guardrails e avaliação por indicador automático de legibilidade.

\section{Guardrails e mitigação de riscos em LLMs}
\label{sec:guardrails-mitigacao-riscos-llms}

Além dos trabalhos voltados à educação de pacientes, mHealth e simplificação textual, também são relevantes os estudos que discutem mecanismos de controle para reduzir riscos no uso de LLMs. Dong et al. ([VERIFICAR\_CITACAO: DONG et al., 2024]) abordam guardrails em um escopo geral, analisando formas de limitar, monitorar e controlar o comportamento de modelos de linguagem em diferentes contextos de aplicação. Embora o trabalho não seja específico da área da saúde, ele contribui para este TCC ao fundamentar a necessidade de mecanismos que restrinjam a geração do modelo e reduzam respostas fora do escopo esperado.

Gangavarapu ([VERIFICAR\_CITACAO: GANGAVARAPU, 2024]), por sua vez, aproxima essa discussão do domínio da saúde ao destacar que aplicações baseadas em IA generativa nesse contexto exigem cuidados adicionais relacionados à factualidade, segurança e confiabilidade das informações produzidas. O trabalho reforça que mecanismos genéricos de controle podem ser insuficientes quando aplicados a cenários sensíveis, como saúde, nos quais erros, omissões ou alterações de sentido podem impactar a compreensão do usuário.

A diferença em relação a este TCC está no nível de aplicação. Enquanto Dong et al. ([VERIFICAR\_CITACAO: DONG et al., 2024]) discutem guardrails para LLMs de forma ampla e Gangavarapu ([VERIFICAR\_CITACAO: GANGAVARAPU, 2024]) trata da segurança de IA generativa em saúde de maneira mais geral, este trabalho considera esses mecanismos no contexto específico da simplificação e adaptação linguística de textos técnicos de saúde. Assim, os guardrails são empregados como apoio para delimitar o uso do modelo, restringir a geração ao conteúdo fornecido e reduzir riscos de criação de informações novas, alterações de sentido ou respostas fora do escopo da tarefa de simplificação.

\section{Comparação dos trabalhos relacionados}
\label{sec:comparacao-trabalhos-relacionados}

A Tabela~\ref{tab:comparacao-trabalhos-relacionados} sintetiza a relação entre os trabalhos analisados e os principais elementos considerados neste TCC. A comparação mostra que os estudos discutidos abordam aspectos importantes para a proposta, como educação de pacientes, letramento em saúde, mHealth, simplificação textual e mitigação de riscos em LLMs. No entanto, esses aspectos aparecem de forma distribuída entre diferentes trabalhos, sem que todos sejam tratados simultaneamente em uma mesma solução aplicada.

\begin{table}[htbp]
  \centering
  \caption{Síntese dos trabalhos relacionados.}
  \label{tab:comparacao-trabalhos-relacionados}
  \small
  \renewcommand{\arraystretch}{1.2}
  \begin{tabularx}{\textwidth}{p{3.0cm} X X}
    \toprule
    \textbf{Trabalho} &
    \textbf{Contribuição para este TCC} &
    \textbf{Limitação em relação a este TCC} \\
    \midrule

    Aydin et al. (2024) &
    Discute o uso de LLMs na educação de pacientes e na adaptação de informações médicas para linguagem mais acessível. &
    Possui escopo amplo e não se concentra especificamente em mHealth, planos de cuidado ou guardrails para simplificação textual. \\

    \addlinespace

    Loughran et al. (2024) &
    Relaciona LLMs, letramento em saúde e mHealth, aproximando-se diretamente do problema investigado neste trabalho. &
    Não enfatiza a adaptação ao perfil do paciente nem a aplicação de guardrails como parte central da solução. \\

    \addlinespace

    Jeblick et al. (2024) &
    Avalia o uso de LLMs para simplificação de textos médicos, evidenciando benefícios e riscos da reescrita automática. &
    Concentra-se em laudos radiológicos e não aborda o contexto mHealth nem a adaptação comunicacional por perfil. \\

    \addlinespace

    Dong et al. (2024) &
    Fundamenta a discussão sobre guardrails como mecanismos de controle para reduzir riscos no uso de LLMs. &
    Trata guardrails em escopo geral, sem foco específico em textos de saúde ou simplificação para pacientes. \\

    \addlinespace

    Gangavarapu (2024) &
    Aproxima a discussão de guardrails do domínio da saúde, destacando riscos de segurança e confiabilidade. &
    Não aborda diretamente simplificação textual, mHealth ou adaptação comunicacional ao perfil do paciente. \\

    \bottomrule
  \end{tabularx}
\end{table}

Aydin et al. ([VERIFICAR\_CITACAO: AYDIN et al., 2024]) oferecem uma visão ampla sobre o uso de LLMs na educação de pacientes, contribuindo para justificar a relevância desses modelos na comunicação em saúde. Loughran et al. ([VERIFICAR\_CITACAO: LOUGHRAN et al., 2024]), por sua vez, aproximam-se mais diretamente deste TCC ao relacionar LLMs, letramento em saúde e aplicações mHealth. Já Jeblick et al. ([VERIFICAR\_CITACAO: JEBLICK et al., 2024]) contribuem especificamente para a discussão sobre simplificação de textos médicos com LLMs, evidenciando tanto o potencial desses modelos quanto os riscos de omissões, distorções e interpretações inadequadas. Por fim, Dong et al. ([VERIFICAR\_CITACAO: DONG et al., 2024]) e Gangavarapu ([VERIFICAR\_CITACAO: GANGAVARAPU, 2024]) fundamentam a importância de mecanismos de controle e mitigação de riscos, especialmente em contextos sensíveis como o da saúde.

A comparação evidencia que a contribuição deste TCC não está em propor isoladamente o uso de LLMs em saúde, nem em desenvolver uma nova métrica de legibilidade ou uma nova plataforma mHealth. A contribuição está na combinação aplicada desses elementos: uma implementação funcional baseada em LLMs para simplificação e adaptação linguística de textos técnicos de saúde, considerando perfil do paciente, guardrails para delimitação da geração e indicadores automáticos de legibilidade e complexidade textual.

Essa combinação delimita uma lacuna prática: como utilizar LLMs para simplificar textos técnicos de saúde de forma adaptada ao perfil do paciente, mantendo a geração restrita ao conteúdo fornecido e reduzindo riscos de alteração de sentido ou criação de informações novas. Assim, este TCC se posiciona entre estudos sobre LLMs na educação de pacientes, letramento em saúde em mHealth, simplificação médica e guardrails, aplicando esses conceitos em uma implementação funcional voltada ao apoio à compreensão de textos de saúde.

% As referências listadas no texto original devem ser cadastradas no arquivo referencias.bib.
% Depois disso, substitua os marcadores [VERIFICAR_CITACAO: ...] por comandos \cite{...}.
%
% Referências a cadastrar:
%
% AYDIN, Serhat; KARABACAK, Mert; VLACHOS, Victoria; MARGETIS, Konstantinos. Large language models in patient education: a scoping review of applications in medicine. Frontiers in Medicine, v. 11, 2024. DOI: 10.3389/fmed.2024.1477898.
%
% DONG, Yi et al. Position: Building Guardrails for Large Language Models Requires Systematic Design. In: Proceedings of the 41st International Conference on Machine Learning. PMLR, v. 235, p. 11375--11394, 2024.
%
% GANGAVARAPU, Ananya. Enhancing Guardrails for Safe and Secure Healthcare AI. arXiv preprint, arXiv:2409.17190, 2024.
%
% JEBLICK, Katharina et al. ChatGPT makes medicine easy to swallow: an exploratory case study on simplified radiology reports. European Radiology, v. 34, p. 2817--2825, 2024. DOI: 10.1007/s00330-023-10213-1.
%
% LOUGHRAN, Elliot; KANE, Madison; WYATT, Tami H.; KERLEY, Alex; LOWE, Sarah; LI, Xueping. Using Large Language Models to Address Health Literacy in mHealth: Case Report. CIN: Computers, Informatics, Nursing, v. 42, n. 10, p. 696--703, 2024. DOI: 10.1097/CIN.0000000000001152.
